\documentclass[12pt]{book}
\usepackage{amsmath}
\usepackage{amssymb}
\usepackage{geometry}
\geometry{margin=1in}

\title{Volume 01: Foundations and Spation Primitives \\ Book 02: Treatise Integration and Foundational Compendium}
\author{James C. Harvey}
\date{December 2025}

\begin{document}
\maketitle

\tableofcontents

\chapter{Treatise Core: The Foundational Compendium}

\section*{Abstract}
This book integrates the SDT treatise corpus into a single narrative line. It establishes the matrix,
flux, medium, and spation as the concrete language of mechanism. The objective is synthesis: a longform
articulation that carries the reader from definitions to consequences with explicit equations embedded
at each step.

\section{Why a Treatise}
SDT is not a collection of disconnected papers; it is a single framework whose equations must be read as
a unified chain. The treatise exists to make that chain explicit. Each derivation is threaded into the next,
so that no statement is free-floating and no formula appears without context or geometry.

\section{The Matrix, Flux, and the Spation}
Space is not empty. The matrix is the continuous collective that supports flux. A spation is the discrete
unit of that matrix; the lattice they form is the medium through which all interaction propagates.
The pressure at a point is not a background constant, but an integral of directional flux reduced by
occlusion:
\begin{equation}
\Pi(\mathbf{r}) = \int_{\Omega} I_{\text{CMB}}(\hat{\mathbf{n}})\left[1 - E(\mathbf{r},\hat{\mathbf{n}})\right]\, d\Omega.
\end{equation}
Here $E$ is occlusion and the integration spans $4\pi$ steradians, establishing that interaction strength
is geometric and directional rather than arbitrary.

\section{Matter as Boundary}
Matter is the bounded geometry that excludes matrix volume. It is not a synonym for displacement, nor a
second substance. It is a boundary condition placed into the matrix that induces flux deficit and creates
gradient-driven movement.

\section{Movement and Now}
Movement is the total class of change in SDT. Shunt dynamics are one facet of movement; circulation,
drift, and flux rearrangement are others. Now is the ordering of change as realized. It is not a new
dimension, but the presentness that accompanies boundary interaction and flux exchange.

\section{The Master Equation as Narrative Spine}
All SDT chapters converge on the master equation for throughput:
\begin{equation}
\dot{E} = P_{\text{CMB}} A_{\text{eff}} \Gamma \kappa (1-\eta).
\end{equation}
Each term is geometric:
\begin{itemize}
\item $P_{\text{CMB}}$ is the universal influx.
\item $A_{\text{eff}}$ is the boundary capture area.
\item $\Gamma$ is circulation ratio.
\item $\kappa$ is curvature.
\item $(1-\eta)$ is coupling efficiency.
\end{itemize}
This equation is the narrative spine of SDT; every phenomenon is a regime of this expression.

\section{Expanded Throughput Derivation}
Write the effective capture area as a geometric projection of boundary surface:
\begin{equation}
A_{\text{eff}} = \int_{\partial \mathcal{B}} \cos\theta(\mathbf{r},\hat{\mathbf{n}})\, dA,
\end{equation}
where $\theta$ is the local incidence angle between flux direction and surface normal. Substituting into
the master equation yields the explicit flux-weighted form:
\begin{equation}
\dot{E} = P_{\text{CMB}} \Gamma \kappa (1-\eta)\int_{\partial \mathcal{B}} \cos\theta\, dA.
\end{equation}
This makes the throughput a direct integral of boundary geometry.

\section{Numerical Substitution Template}
Using canonical constants, the master equation can be written as
\begin{equation}
\dot{E} = (2.036 \times 10^{-2})\ A_{\text{eff}}\ \Gamma\ \kappa\ (1-\eta)\ \text{W},
\end{equation}
where $A_{\text{eff}}$ is in $\text{m}^2$ and $\kappa$ in $\text{m}^{-1}$. This expression is the
numerical template applied in benchmarks.

\section{Derivation Tree}
The master equation unfolds into a derivation tree:
\begin{align}
\Pi(\mathbf{r}) &\rightarrow \nabla \Pi(\mathbf{r}) \rightarrow \mathbf{F},\\
\mathbf{F} &\rightarrow \text{movement} \rightarrow \text{structure},\\
\Gamma,\kappa,\eta &\rightarrow \mu,\, E_{\text{bind}},\, \Delta E_{\text{spec}}.
\end{align}
Each arrow is a geometric constraint, not a hypothesis.

\section{Data Compendium Integration}
The treatise integrates derived tables and numeric compendia as direct projections of the master
equation. For a given boundary class $\mathcal{B}$, each entry records
\begin{equation}
\left\{A_{\text{eff}}, \Gamma, \kappa, \eta, \dot{E}\right\},
\end{equation}
so that every numeric value can be traced to geometric inputs.

\section{Calculator Suite Consistency}
The computational suite implements the same derivation chain:
\begin{equation}
\dot{E}(\mathcal{B}) \rightarrow E_{\text{bind}} \rightarrow \Delta E_{\text{spec}} \rightarrow \text{observable}.
\end{equation}
Each step is documented and reproducible without parameter tuning.

\section{Constants as Geometric Anchors}
The constants are anchors, not arbitrary fittings:
\begin{align}
P_{\text{CMB}} &= 2.036 \times 10^{-2}\ \text{Pa},\\
K_{\text{bulk}} &= 4.6 \times 10^{113}\ \text{Pa},\\
c &= 2.99792458 \times 10^8\ \text{m/s}.
\end{align}

\section{Dimensional Closure}
Every formula closes dimensionally:
\begin{align}
[\dot{E}] &= \text{W} = \text{kg}\cdot \text{m}^2\cdot \text{s}^{-3},\\
[P] &= \text{Pa} = \text{kg}\cdot \text{m}^{-1}\cdot \text{s}^{-2},\\
[\kappa] &= \text{m}^{-1},\quad [\Gamma]=1,\quad [\eta]=1.
\end{align}

\chapter{CMB as the Universal Source}

\section*{Abstract}
The CMB is not a background relic; it is the active flux reservoir. This chapter derives how isotropic
flux becomes local pressure and how occlusion converts that pressure into interaction.

\section{Pressure from Flux}
The pressure corresponding to isotropic flux is
\begin{equation}
P_{\text{CMB}} = 4\pi I_{\text{CMB}}.
\end{equation}
Local pressure is the occlusion-modulated integral in Equation (1.1). Therefore the existence of matter
is the existence of occlusion, and the existence of occlusion is the existence of gradients.

\section{Radial Deficit Field}
A bounded displacement induces a radial pressure deficit. In the far field, the pressure profile takes
the form
\begin{equation}
\Pi(r) = \Pi_0 - \frac{\beta \rho_s}{r},
\end{equation}
where $\Pi_0$ is the ambient matrix pressure, $\rho_s$ is the matrix density, and $\beta$ is the
geometric potential parameter derived from displacement volume.

\section{Acceleration from Pressure Gradient}
The gradient of the deficit produces acceleration:
\begin{equation}
\frac{d\Pi}{dr} = \frac{\beta \rho_s}{r^2},\quad
a(r) = -\frac{1}{\rho_s}\frac{d\Pi}{dr} = -\frac{\beta}{r^2}.
\end{equation}
The inverse-square behavior is therefore a pressure-gradient consequence, not a separate postulate.

\section{Geometric Potential Parameter}
The parameter $\beta$ is defined by displacement geometry:
\begin{equation}
\beta = \frac{\kappa V_{\text{disp}} K_{\text{bulk}}}{4\pi \rho_s} = \frac{\kappa V_{\text{disp}} c^2}{4\pi}.
\end{equation}
This replaces mass and gravitational constant with a single geometric displacement parameter.

\section{Screening and Force Hierarchy}
Composite boundaries screen their own displacement fields. Define the effective displacement fraction
$\xi$ by
\begin{equation}
\xi = \frac{V_{\text{disp,eff}}}{V_{\text{disp,total}}},
\end{equation}
with $\xi \ll 1$ for large composite bodies. This screening accounts for the disparity between
short-range coupling and long-range coupling within a single pressure field.

\section{Refractive Index Gradient}
Pressure gradients alter propagation as a refractive index field:
\begin{equation}
n(r) = 1 - \frac{\beta}{c^2 r}.
\end{equation}
This yields deflection of flux paths in the presence of large displacement boundaries.

\section{Orbital Precession Term}
Higher-order pressure terms yield an effective correction
\begin{equation}
\Phi_{\text{eff}}(r) = -\frac{\beta}{r} - \frac{3\beta^2}{2c^2 r^2},
\end{equation}
which produces a measurable precession in long-term orbital motion without invoking external curvature
postulates.

\section{Hierarchy of Regimes}
Small boundaries with high curvature intercept more flux and yield strong coupling; large boundaries with
low curvature yield weak coupling. The hierarchy of forces is therefore the hierarchy of geometry.
\begin{equation}
F(\kappa_1,\kappa_2,r) = \frac{\kappa_1 \kappa_2}{4\pi K_{\text{bulk}} r^2}\,\mathcal{C}(\mathbf{r},\hat{\mathbf{n}}).
\end{equation}
No new forces are introduced—only new regimes of the same equation.

\section{Flux Balance and Equilibrium}
Equilibrium is the condition of zero net gradient:
\begin{equation}
\nabla \Pi(\mathbf{r}) = 0.
\end{equation}
Stable structures are those that minimize net gradient while preserving internal circulation.

\chapter{Historical Foundations and Consolidations}

\section*{Abstract}
The historical archive is not discarded; it is consolidated. This chapter describes how prior phases map
into the current SDT grammar and where older terminology is translated into current terms.

\section{Translation Rules}
\begin{itemize}
\item ``Spation'' denotes a discrete unit; the ``matrix'' denotes the collective medium.
\item ``Flux'' refers to directional propagation; ``medium'' refers to the collective substrate.
\item ``Movement'' is the class; ``shunt'' is a mechanism within the class.
\end{itemize}

\section{Continuity of Equations}
All archived derivations that express pressure gradients and curvature coupling are preserved, but the
notation is unified so that every formula aligns with the master equation and occlusion integral.

\section{Phase Mapping}
Historical phases are re-indexed against the four primitives and the master equation. A prior result is
canonical only if it can be written in the form
\begin{equation}
\mathcal{R}_i: \dot{E} = P_{\text{CMB}} A_{\text{eff}} \Gamma \kappa (1-\eta),
\end{equation}
with explicit definitions for each geometric factor.

\section{Canonicalization}
Canonicalization means that every statement is mapped to one of the four primitives. If a formula does
not transparently reduce to space, matter, movement, and now, it is re-expressed until it does.

\chapter{A Compendium of Constants}

\section*{Abstract}
SDT constants are not arbitrary; they are the numeric anchors of geometry. This chapter lists and uses
them in-line with formulae to maintain dimensional clarity.

\section{Core Constants}
\begin{align}
P_{\text{CMB}} &= 2.036 \times 10^{-2}\ \text{Pa},\\
K_{\text{bulk}} &= 4.6 \times 10^{113}\ \text{Pa},\\
c &= 2.99792458 \times 10^8\ \text{m/s}.
\end{align}
These constants are used in all subsequent derivations.

\section{Derived Measures}
From $K_{\text{bulk}}$ and $c$, the effective matrix density is
\begin{equation}
\rho_s = \frac{K_{\text{bulk}}}{c^2}.
\end{equation}
This quantity anchors all pressure-gradient acceleration expressions in the treatise.

\section{Derived Constants}
Derived constants emerge from geometry:
\begin{equation}
G = \frac{c^4}{4\pi K_{\text{bulk}}^2}.
\end{equation}
These expressions are not assumptions; they are consequences.

\section{Explicit Substitution}
Substituting the canonical values gives the numerical construction:
\begin{equation}
G = \frac{(2.99792458 \times 10^8)^4}{4\pi (4.6 \times 10^{113})^2}.
\end{equation}
The expression shows that $G$ is a consequence of matrix stiffness and wave speed.

\chapter{The Narrative of Derivation}

\section*{Abstract}
This chapter is a continuous derivation narrative, meant to be read as a single mathematical argument.
It moves from flux to pressure, from pressure to gradient, from gradient to movement, and from movement
to structure.

\section{Flux to Pressure}
\begin{equation}
\Pi(\mathbf{r}) = \int_{\Omega} I_{\text{CMB}}(\hat{\mathbf{n}})\left[1 - E(\mathbf{r},\hat{\mathbf{n}})\right]\, d\Omega.
\end{equation}

\section{Pressure to Movement}
\begin{equation}
\mathbf{F} = -\nabla \Pi(\mathbf{r}),
\end{equation}
and movement follows as the reconfiguration of boundaries under the gradient.

\section{Movement to Structure}
The stable configuration is the one that minimizes net gradients while maximizing consistent flux
throughput, yielding the observed lattices and boundary geometries.

\section{Structure to Spectrum}
Spectral spacing follows from structured circulation:
\begin{equation}
\Delta E \propto \Gamma \kappa (1-\eta),
\end{equation}
so line spacing is the geometric signature of boundary topology.

\chapter{Methodological Discipline}

\section*{Abstract}
Method in SDT is geometric discipline. Every claim is anchored by a boundary, every equation by a
dimension, every derivation by a primitive.

\section{No Fitted Parameters}
Parameters are derived, not tuned. If a quantity cannot be derived, it is not canonical.

\section{Validation Path}
Validation proceeds by tracing symbolic expressions into numeric predictions and comparing them with
observed values without adjustment.

\section{Dimensional Audit}
Every derivation is verified by unit propagation:
\begin{equation}
[P][A][\Gamma][\kappa] \rightarrow \text{W},
\end{equation}
ensuring the master equation remains dimensionally closed at every stage.

\chapter{Closing Synthesis}

\section*{Abstract}
The treatise is not an appendix; it is the spine. This closing chapter ties the narrative thread into a
single statement.

\section{One Sentence}
Space is the matrix, matter is boundary, movement is flux-driven change, now is the ordering of change—
and every equation in SDT is the algebra of that sentence.

\end{document}
